% ============================================================================
% THE CORE FIFTY --- program-wide triage tier, shared by all four volumes.
% Canonical source: kernel/core50.tex; copied into each volume by the build.
% Constraints: no \ref (this file is shared), no macro definitions, no tikz,
% no interviewq. Cross-volume pointers are textual: [DL n], [NLP n], [SR n],
% [CML n].
% ============================================================================

\chapter*{The Core Fifty}
\addcontentsline{toc}{chapter}{The Core Fifty}

\section*{How to Use This Tier}
\addcontentsline{toc}{section}{How to Use This Tier}

\begin{tldr}
Fifty things a Staff-level candidate cannot skip, drawn from all four volumes. Each entry gives the substance you must be able to state, where the full treatment lives, and the specific wrong answer that costs the signal. This is a checklist for the last week --- not a substitute for the chapters, and not a reading list.
\end{tldr}

The program holds 543 questions across four volumes. You will not master all of them, and you do not need to: interview loops are not uniform samples. A small set of topics shows up in nearly every loop, carries the weight of the answers around it, and separates a strong Staff candidate from a good Senior one. That set is what this tier names.

\textbf{The selection rule.} Items earned a place by four criteria, in priority order. \emph{Frequency}: it shows up in most loops at this level. \emph{Load-bearing}: other answers depend on it, so not knowing it collapses a whole area rather than costing one question. \emph{Discriminating}: it is where Senior and Staff answers visibly diverge --- usually because the Staff answer carries a number, a mechanism, or a decision rule where the Senior answer carries a name. \emph{Coverage}: the fifty must span breadth, depth, coding, and design, and span all four volumes.

\textbf{The ninety-second rule.} An item you cannot explain out loud, from memory, in ninety seconds, is not yet yours. Reading an entry and nodding is recognition, not retrieval, and interviews test retrieval under follow-up pressure. Speak each entry aloud. Record yourself for the ones you doubt. The gap between what you can recognize and what you can say is the entire subject of this tier.

\textbf{How to run it.} Five items a day for ten days, spoken, marked red / amber / green. Green items are done --- do not reread them. Amber items get one more pass later in the week. Only red items earn chapter time, and then you read the pointer given here rather than the whole chapter. In the final forty-eight hours, run the checklist table at the end of this chapter and speak only the reds.

\begin{keyinsight}[The Core Fifty Is Not the Kernel]
The kernel is the set of principles you reason \emph{from} --- the handful of ideas that let you re-derive a forgotten detail in the room. The Core Fifty is specific material you must have \emph{ready}: formulas, numbers, mechanisms, and decision rules that should arrive without derivation. Many entries here instantiate a kernel principle; that is expected. The test is different, though. For a kernel page the test is ``can I rebuild this from scratch?'' For a Core Fifty entry the test is ``does this come out of my mouth in ninety seconds, with the number in it?''
\end{keyinsight}

\begin{warningbox}
Do not use this tier as a reading list, and do not use it to decide what to skip in your own domain. If you interview for a search team, all of Volume III is in scope regardless of what these ten retrieval entries cover; if you interview for a frontier-lab pretraining team, [DL 25] and [DL 26] are in scope well beyond the four entries drawn from them. The Core Fifty is the floor, not the ceiling --- it is what you must have when the loop wanders outside your specialty, which it always does.
\end{warningbox}

% ============================================================================
\section*{I. Foundations}
\addcontentsline{toc}{section}{I. Foundations}
% ============================================================================

\subsection*{1. The softmax cross-entropy gradient}

\textbf{What you must be able to say.} For $L = -\log \mathrm{softmax}(z)_y$, the gradient with respect to the logits is $\partial L / \partial z = p - y$: predicted probability minus the one-hot target, bounded in $[-1, 1]$ per coordinate. Softmax alone has Jacobian $\mathrm{diag}(p) - pp^{\top}$, which saturates; composing it with cross-entropy cancels that factor exactly, which is why the pair never suffers vanishing gradients at the output layer and why you always fuse the two (log-sum-exp with the max subtracted) rather than computing softmax then log.

\textbf{Where it lives.} [DL 1], ``Derive the gradient of softmax cross-entropy loss with respect to the logits.'' The numerical-stability tricks are in the same chapter; the loss taxonomy that hangs off this result is [DL 3].

\textbf{The failure.} Quoting $p - y$ without being able to produce the two-line derivation, or --- worse --- deriving the softmax Jacobian and stopping there, missing that the cancellation with cross-entropy is the whole point. A close second: not knowing why MSE on logits is a bad classification loss (tiny gradients exactly when the model is confidently wrong).

\subsection*{2. Perplexity, cross-entropy, and bits}

\textbf{What you must be able to say.} Perplexity is the exponentiated per-token cross-entropy: $\mathrm{PP} = \exp(-\frac{1}{N}\sum_i \log p(w_i \mid w_{<i})) = e^{H_{\text{nats}}} = 2^{H_{\text{bits}}}$, and it reads as an effective branching factor --- perplexity 50 means the model is as uncertain as a uniform choice among 50 tokens. Cross-entropy decomposes as $H(p) + \mathrm{KL}(p \Vert q)$, so the loss has an irreducible floor and a loss delta is literally bits saved. The limitation that matters: perplexity is defined per token, so it is not comparable across tokenizers --- normalize to bits per character or bits per byte before you compare two models.

\textbf{Where it lives.} [NLP 10], ``What is perplexity? How does it relate to cross-entropy? What are its limitations?'' The same question is canon in [NLP 2]; the entropy and KL machinery is [DL 1].

\textbf{The failure.} Comparing perplexities across models with different vocabularies, and not volunteering the bits-per-byte fix. The other failure is treating perplexity as a quality metric: a model can win on perplexity and lose on every generation eval, because perplexity scores the whole distribution while users experience the argmax.

\subsection*{3. Double descent and why over-parameterization generalizes}

\textbf{What you must be able to say.} Test error follows the classical U-curve up to the interpolation threshold (roughly, parameters enough to fit the training set exactly), spikes there, then descends again as capacity grows past it. Among the infinitely many zero-training-loss solutions, gradient descent's implicit bias selects low-norm, flat ones, and larger models have more such basins --- so bigger generalizes better, not worse. The bias-variance tradeoff is a correct description of the underparameterized regime only. The modern-LLM corollary: frontier pretraining runs see most tokens once, so train loss $\approx$ test loss and classical overfitting barely applies.

\textbf{Where it lives.} [DL 2], ``Explain double descent'' and ``Explain the bias-variance tradeoff. Does it apply to deep learning?'' Grokking, the lottery ticket, and implicit regularization are in the same chapter.

\textbf{The failure.} Reciting the bias-variance tradeoff as though it settles the question, then being unable to reconcile it with the fact that every frontier lab's answer to overfitting is a bigger model. Also: claiming 100\% training accuracy is proof of overfitting.

\subsection*{4. The 6ND rule and Chinchilla allocation}

\textbf{What you must be able to say.} Training compute is $C \approx 6ND$ FLOPs for $N$ parameters and $D$ tokens (2 FLOPs per parameter per token forward, roughly twice that backward). For a fixed $C$, loss is minimized with $D \approx 20N$, so both scale as $C^{0.5}$. Worked: at $C = 10^{22}$, $C = 120 N^2$ gives $N \approx 9\times 10^{9}$ parameters on $D \approx 180$B tokens. Then the Staff caveat: compute-optimal is the wrong target when you will serve the model, because total cost is training plus inference times queries --- which is why Llama-class models train far past 20:1 (Llama 3 8B on 15T tokens is roughly 1900:1).

\textbf{Where it lives.} [DL 2], ``Estimate the compute-optimal model size for a training budget of $10^{22}$ FLOPs using Chinchilla scaling laws'' and the fixed-budget size-versus-data question. Scaling laws as an LLM topic are also [NLP 7].

\textbf{The failure.} Knowing ``20 tokens per parameter'' as a slogan but not being able to do the algebra from $C = 6ND$ to $N \propto C^{0.5}$, and not raising the inference-cost caveat unprompted. Quoting Kaplan's exponents as current is a second-order tell.

% ============================================================================
\section*{II. Architectures}
\addcontentsline{toc}{section}{II. Architectures}
% ============================================================================

\subsection*{5. Self-attention, the $\sqrt{d_k}$ scaling, and multiple heads}

\textbf{What you must be able to say.} $\mathrm{Attention}(Q,K,V) = \mathrm{softmax}(QK^{\top}/\sqrt{d_k})V$ is a soft dictionary lookup: queries score keys by dot product, and the output is the value average under those scores. The scaling is a variance argument, not a convention: with unit-variance i.i.d.\ entries, $q \cdot k$ is a sum of $d_k$ products and so has variance $d_k$ and standard deviation $\sqrt{d_k}$; dividing restores unit variance and keeps the softmax out of its saturated regime, where gradients vanish. Multiple heads are $h$ parallel lookups in $d/h$-dimensional subspaces at identical total FLOPs --- one head can only attend to one weighted average per position, so it must blend distinct relations (syntactic, positional, coreferent) into a single distribution.

\textbf{Where it lives.} [NLP 5], ``Explain self-attention step by step,'' ``Why do we scale by $d_k$? Derive the variance argument,'' and ``Why do we need multiple attention heads?'' [DL 5] carries the same block as a cost model.

\textbf{The failure.} ``We divide by $\sqrt{d_k}$ to normalize'' without the variance calculation --- interviewers push here precisely because it separates memorization from understanding. On heads, the failure is saying ``more capacity'' rather than naming what a single softmax-weighted average cannot represent.

\subsection*{6. Attention versus FFN cost, and the crossover}

\textbf{What you must be able to say.} Per layer per token, attention costs about $8d^2 + 4nd$ FLOPs (QKV and output projections, plus scores and value-weighting) while a SwiGLU FFN costs about $16d^2$, so the ratio is $0.5 + n/(4d)$. Attention only overtakes the FFN at $n \approx 2d$ --- around 8K tokens at $d = 4096$. That is the number that kills the reflex answer: at ordinary context lengths a transformer is FFN-bound and quadratic attention is not your problem; at 128K it is. Complexity is $O(n^2 d)$ for attention against $O(n d^2)$ for the FFN, which is the whole reason long-context work targets attention specifically.

\textbf{Where it lives.} [NLP 5], ``What is the computational complexity of self-attention? Why is it a problem?'' The worked ratio and the crossover arithmetic are [DL 5], ``Estimate the ratio of attention FLOPs to FFN FLOPs at different sequence lengths.''

\textbf{The failure.} Saying ``attention is quadratic, that's the bottleneck'' at every sequence length. At 2K context on a 7B model the FFN dominates by a wide margin, and a candidate who cannot say where the crossover is has not costed a transformer.

\subsection*{7. Encoder-only, decoder-only, encoder--decoder}

\textbf{What you must be able to say.} Encoder-only (BERT, masked LM) sees bidirectional context and produces representations --- best for classification, retrieval embeddings, and token labelling, but it cannot generate. Decoder-only (GPT, causal LM) predicts the next token with a causal mask, which makes every token a training signal, makes KV caching possible, and makes one model serve every task through prompting. Encoder--decoder (T5) keeps a bidirectional encoder for a fixed input plus a decoder for output, and still wins where the input is fixed and the output is a transformation of it. Decoder-only won because the objective is denser per token, scales cleanly, and unifies tasks --- not because it is architecturally superior.

\textbf{Where it lives.} [NLP 5], ``Compare encoder-only, decoder-only, and encoder-decoder transformers. Why did decoder-only dominate the LLM era?'' The objective-level comparison (MLM against CLM) is [DL 8], ``BERT vs GPT pre-training objectives.''

\textbf{The failure.} Answering with a table of model names and no mechanism. The Staff answer names the training-signal density argument (MLM supervises the 15\% of positions that are masked; CLM supervises every position) and still concedes that an encoder is the right choice for a 50-category classifier at high QPS.

\subsection*{8. Vanishing gradients, LSTMs, and why transformers won}

\textbf{What you must be able to say.} Backpropagation through time multiplies Jacobians across timesteps, so gradients decay or explode geometrically in sequence length; the LSTM's cell state gives an \emph{additive} path ($c_t = f_t \odot c_{t-1} + i_t \odot \tilde{c}_t$) whose gradient is gated rather than repeatedly multiplied by a weight matrix, which extends the usable horizon to hundreds of steps rather than thousands. Transformers replaced LSTMs for two reasons, and you must give both: constant path length between any two positions (no decay at all), and full parallelism across the sequence at training time, where an RNN is inherently sequential. The cost is $O(n^2)$ attention instead of $O(n)$ recurrence.

\textbf{Where it lives.} [NLP 4], ``Explain the vanishing gradient problem in RNNs. Why does LSTM solve it?'' and ``Why did transformers replace LSTMs?''

\textbf{The failure.} Giving only the parallelism argument. Parallelism explains why transformers train faster; path length explains why they model long-range dependencies better. A candidate who names only one has half the answer --- and will be stuck when asked why Mamba, which is sequential at inference, is competitive.

\subsection*{9. The modern block: pre-norm, RMSNorm, SwiGLU}

\textbf{What you must be able to say.} Pre-norm puts the normalization inside the residual branch, leaving the residual stream an unobstructed identity path, so gradients reach layer 1 without passing through a normalization at every level --- that is why 100-layer models train without warmup gymnastics, at the cost of slightly weaker representations than post-norm. RMSNorm drops mean-centering and keeps only the scale, saving roughly 10--15\% of normalization time with no quality cost, because pre-norm already preserves the stream's mean. SwiGLU replaces the FFN's single nonlinearity with a gated pair, $(\mathrm{Swish}(xW) \odot xV)W_2$, and to keep parameters constant the hidden width shrinks to $\tfrac{8}{3}d$ --- 11008 at $d = 4096$ for Llama-2-7B. The FFN is where about two-thirds of the parameters live (4.3B of 6.7B).

\textbf{Where it lives.} [DL 4], ``Why GELU over ReLU? Why SwiGLU over GELU?'', ``RMSNorm vs LayerNorm,'' and the pre-norm versus post-norm gradient-flow question. The architecture-evolution framing is [DL 5], ``What changed between GPT-2 and LLaMA?''

\textbf{The failure.} Listing the changes without the reason for each. Every item in the modern block is a fix for a named failure --- and the tell of a strong candidate is being able to say that post-norm is not simply worse, only harder to train deep.

\subsection*{10. RoPE and context extension}

\textbf{What you must be able to say.} RoPE rotates query and key vectors by an angle proportional to absolute position, so the dot product depends only on the \emph{relative} offset --- position enters multiplicatively inside the score rather than as an added embedding. That is why it composes unchanged with FlashAttention and the KV cache (each key is rotated once, at its own position, before it is written). Extending context means manipulating the base frequency: position interpolation squeezes positions into the trained range, NTK-aware scaling and YaRN adjust frequencies per band, and all of them need a modest continued-pretraining pass to recover quality. Degradation past the trained length is not a bug in attention; it is out-of-distribution rotation angles.

\textbf{Where it lives.} [DL 5] for RoPE at working depth, ``How does RoPE encode position information?'' The context-extension toolbox and its costs are [DL 21], ``Your 8K-trained model must serve 64K contexts next month.'' The scheme-by-scheme comparison against sinusoidal and learned embeddings is [NLP 5].

\textbf{The failure.} Describing RoPE as ``relative position embeddings'' without the rotation mechanism, and then having no answer for why it extends to longer context better than learned absolute positions. The second failure is proposing to change the context length as a config flag with no continued training and no evaluation at the new length.

\subsection*{11. Mixture of experts}

\textbf{What you must be able to say.} An MoE replaces the FFN with $E$ experts and a router that sends each token to its top-$k$, so parameters scale with $E$ while FLOPs per token scale with $k$ --- DeepSeek-V3 is 671B total with 37B active. The fine-grained trend (256 small experts, top-8, rather than 8 large ones, top-2) buys combinatorial specialization at fixed active-parameter budget. The load-balancing problem is the whole engineering story: without pressure, the router collapses onto a few experts, so you add an auxiliary balance loss, a capacity factor with token dropping, or DeepSeek-style bias adjustment outside the gradient path. MoE is memory-hungry and communication-hungry (two all-to-alls per layer per pass), so it wins on training-compute efficiency and loses when serving memory or small-batch decode is the constraint.

\textbf{Where it lives.} [DL 14], ``When would you choose MoE over dense scaling?'' and ``MoE routing collapse: what is it, how do you detect it, and how do you fix it?'' The routing and load-balancing conceptual version is [NLP 7]; expert parallelism is [DL 15].

\textbf{The failure.} Selling MoE as free capacity. The Staff answer prices it: you still hold every expert in memory, you pay all-to-all on every layer, an imbalanced route makes one GPU the critical path while the rest idle, and at low batch each expert's GEMM falls below the hardware's compute/bandwidth ridge.

\subsection*{12. Tokenization: BPE and the tokenizer tax}

\textbf{What you must be able to say.} BPE starts from bytes or characters and repeatedly merges the most frequent adjacent pair, recording the merge list; encoding replays those merges in order. Because the base alphabet is bytes, there is no out-of-vocabulary case --- unknown text degrades to more tokens, not to an UNK. The consequences you must carry: token count, not character count, is what you pay for and what fills the context window; a tokenizer trained on English spends two to four times as many tokens per unit of meaning on other languages and on code, which is a latency, cost, and quality tax on exactly the users you did not train for; and vocabulary size trades embedding and softmax cost against sequence length.

\textbf{Where it lives.} [NLP 12], ``Explain BPE tokenization. How is it trained? How does it handle OOV words?'' and the tokenizer-tax question in the same chapter. The implementation drill is [DL 28]; the interaction with embedding quality is [DL 6]; the vocabulary-size decision for a new model family is [DL 25].

\textbf{The failure.} Describing BPE as ``splitting on subwords'' with no merge-training procedure, and not connecting tokenization to anything downstream. A Staff candidate reaches for tokenization when diagnosing why a multilingual product is slow and expensive, or why the model cannot do arithmetic on long numbers.

% ============================================================================
\section*{III. Training and Post-Training}
\addcontentsline{toc}{section}{III. Training and Post-Training}
% ============================================================================

\subsection*{13. The memory ledger: where parameters and bytes go}

\textbf{What you must be able to say.} Mixed-precision AdamW training costs roughly 16--18 bytes per parameter before activations: BF16 weights (2) plus gradients (2--4) plus an FP32 master copy (4) plus two FP32 Adam moments (8). A 7B model is therefore about 112--126\,GB and does not fit one 80\,GB GPU; inference is 2 bytes per parameter (0.5 at INT4) plus the KV cache. Activations scale as batch $\times$ sequence $\times$ hidden $\times$ layers and are traded against roughly 30\% extra compute by gradient checkpointing. You should also be able to place the parameters: a Llama-class model is about $12 L d^2$, so 32 layers at $d = 4096$ gives 6.4B, with the FFN holding two-thirds and embeddings the rest.

\textbf{Where it lives.} [DL 15], ``How much GPU memory is needed to fine-tune a 13B parameter model with LoRA vs.\ full fine-tuning?'' [DL 5] for ``Estimate the memory required to train a 7B parameter model'' and [DL 9] for ``Roughly where do those 7B parameters live?''

\textbf{The failure.} Answering a memory question with parameter count times precision and forgetting optimizer state --- which is half the bill. The second failure is not knowing what ZeRO stages 1, 2, and 3 shard (optimizer state, then gradients, then parameters) when asked how to make it fit.

\subsection*{14. Parallelism: DP, TP, PP, ZeRO, EP}

\textbf{What you must be able to say.} Four axes, chosen by what they cost in communication. Data parallelism replicates the model and all-reduces gradients --- cheapest, but requires the model state to fit, which ZeRO/FSDP fixes by sharding optimizer state, gradients, and parameters across the data-parallel group. Tensor parallelism splits individual matrices and communicates twice per layer, so it must stay inside a node on NVLink. Pipeline parallelism splits layers across nodes and communicates only activations at stage boundaries, so it tolerates slower links but introduces bubbles that micro-batching only partly hides. Expert parallelism adds two all-to-alls per MoE layer. The canonical 70B-on-256-GPU layout: TP 8 within the node, PP across nodes, DP with ZeRO-1 on top, plus activation checkpointing.

\textbf{Where it lives.} [DL 15], ``How would you scale training of a 70B parameter model across 256 GPUs?'' and ``Data parallelism vs.\ tensor parallelism vs.\ pipeline parallelism --- what is your decision framework?'' The link-characteristics mapping is [DL 26].

\textbf{The failure.} Naming the three strategies without mapping them onto the network. The reason TP stays inside a node and PP crosses nodes is bandwidth, and a candidate who cannot say that has memorized a taxonomy rather than reasoned about a cluster.

\subsection*{15. Loss spikes and NaNs}

\textbf{What you must be able to say.} Order the diagnosis by prior probability and name the cheapest discriminating check for each. Data first: a corrupt or pathological shard, checked by replaying the exact batch. Numerics second: FP16 overflow in attention or the loss scaler (BF16's wider exponent is why it is the default), checked by logging gradient and activation norms per layer. Third, and the one that separates candidates: attention logits growing into the hundreds over training, which saturates the softmax and destabilizes the run --- the fixes are QK-norm on queries and keys and a z-loss term $\lambda \log^2 Z$ that keeps the output softmax normalizer near one. Then the cheap knobs: lower peak LR, longer warmup, $\beta_2 = 0.95$, tighter gradient clipping. Operationally: skip the offending step, or rewind to the last checkpoint and skip the batch.

\textbf{Where it lives.} [DL 4], ``Training loss is NaN after a few hundred steps'' and ``Your 70B-parameter run in BF16 starts loss-spiking at 100B tokens.'' The full differential diagnosis at scale is [DL 15], ``Your 70B pretraining run diverges at 200B tokens.''

\textbf{The failure.} An unordered list of everything that can go wrong. The signal here is triage: which hypothesis first, what one measurement distinguishes it, and what it costs to run that measurement on a live billion-dollar job.

\subsection*{16. The pretraining data pipeline}

\textbf{What you must be able to say.} The order is filter, then deduplicate, then decontaminate, then design the mixture --- and the order matters, because deduplicating before filtering wastes the expensive step on junk and decontaminating before deduplication leaves near-duplicate benchmark leakage behind. Filtering is a cascade: language ID, heuristic quality rules, then a learned quality classifier. Deduplication is MinHash/LSH at document level plus exact substring removal; it reliably improves models, but over-aggressive deduplication strips legitimately repeated high-value text. Decontamination by $n$-gram overlap misses paraphrased and translated contamination, so you also hold out fresh evaluations. Mixture weights (code, web, math) and the mid-training anneal are set by small ablation runs against a fixed eval suite, not by argument.

\textbf{Where it lives.} [DL 25], ``Walk me through what happens between `we have a Common Crawl snapshot' and the first training step,'' plus the mixture, decontamination, and mid-training questions in the same chapter. The classical filtering instruments (fastText, KenLM, MinHash) are [NLP 1].

\textbf{The failure.} Treating data as a solved preprocessing step. At Staff level the expected shape of the answer is a program: an ablation cadence that gives every ingredient an evidence trail before a nine-figure run commits to it.

\subsection*{17. LoRA: the count, the ratio, the reason}

\textbf{What you must be able to say.} LoRA freezes $W$ and learns $W + \frac{\alpha}{r}BA$ with $B$ and $A$ of rank $r$. The count is $L \times M \times 2 \times d \times r$: for Llama-7B at $r = 16$ on the four attention projections, $32 \times 4 \times 2 \times 4096 \times 16 = 16.8$M, about 0.25\% of the base; adding the FFN matrices roughly doubles it to 0.6\%. Only the ratio $\alpha/r$ matters --- it is the scale on the update, entangled with learning rate, so fix the ratio (2 is the common default) and tune the LR. Because vanilla LoRA divides by $r$ rather than $\sqrt{r}$, raising rank at fixed $\alpha$ \emph{shrinks} the effective update, which is why a naive $r = 64$ run can look no better than $r = 16$. It works because fine-tuning updates have low intrinsic rank: adaptation moves the model within a small subspace, it does not rebuild it.

\textbf{Where it lives.} [DL 16], ``Estimate the number of trainable parameters for LoRA with rank 16 on a 7B model'' and the rank/$\alpha$/target-module defence question. The adaptation-lifecycle framing is [NLP 6].

\textbf{The failure.} Treating $\alpha$ as a second capacity knob independent of $r$, and not knowing that the memory win comes from optimizer state (no Adam moments for the frozen base), not from the adapter parameters themselves.

\subsection*{18. Catastrophic forgetting}

\textbf{What you must be able to say.} Fine-tuning on a narrow distribution moves weights toward it and away from everything else; the model gets better at your task and worse at the general capabilities you did not measure. The mitigations, in the order you would try them: fewer epochs and a lower learning rate (most ``forgetting'' is over-training); mix 5--30\% general or replay data into the fine-tuning set; prefer a parameter-efficient method with a small $\alpha/r$ so the frozen base dominates; and --- non-negotiable --- run a general-capability eval suite as a guardrail on every checkpoint, not just the task metric. If it has already happened, the honest fix is to retrain with the guardrail in the loop, not to patch.

\textbf{Where it lives.} [NLP 6], ``What is catastrophic forgetting? How do you mitigate it during fine-tuning?'' The LoRA-specific version is [DL 16], ``Your LoRA fine-tuned model performs well on your test set but catastrophically forgets general capabilities.''

\textbf{The failure.} Naming EWC and regularization-based methods from the continual-learning literature while missing the two things that actually work in practice: train less, and mix in replay data. The second failure is not having noticed the problem at all, which is a monitoring failure and reads as such.

\subsection*{19. The RLHF pipeline, the KL anchor, and reward hacking}

\textbf{What you must be able to say.} Three stages: SFT on demonstrations; a reward model trained on preference pairs under a Bradley--Terry likelihood; then policy optimization (PPO) against $r(x,y) - \beta\,\mathrm{KL}(\pi_\theta \Vert \pi_{\text{ref}})$. Stage three holds four models in memory --- policy, frozen reference, reward model, critic --- roughly 56\,GB in FP16 at 7B for parameters alone. The KL term is the answer to the question ``why doesn't it just maximize reward'': the reward model is a proxy that is only accurate near the distribution it was trained on, so unbounded optimization walks off that distribution and finds its blind spots. That is reward hacking, and its concrete forms are verbosity, sycophancy, and format tricks. Detect it by reward climbing while KL climbs and human preference stalls; defend with a KL budget, RM ensembles, early stopping, and fresh held-out human evals.

\textbf{Where it lives.} [DL 17], ``Walk me through the RLHF pipeline,'' ``Why does the KL penalty matter in RLHF? What happens without it?'', and ``What is reward hacking? Give 3 concrete examples.'' The production alignment pipeline and cluster layout are [NLP 8].

\textbf{The failure.} Describing the KL term as a regularizer without saying what it regularizes \emph{toward} and why that matters. Reward hacking is not a curiosity: it is the predictable consequence of optimizing a learned proxy, and the candidate who frames it that way is reasoning, not reciting.

\subsection*{20. DPO}

\textbf{What you must be able to say.} DPO removes the reward model by inverting the KL-constrained RLHF optimum. The optimal policy is $\pi^{*}(y \mid x) \propto \pi_{\text{ref}}(y \mid x)\exp(r(x,y)/\beta)$; solve for the reward, $r(x,y) = \beta \log \frac{\pi^{*}(y\mid x)}{\pi_{\text{ref}}(y \mid x)} + \beta \log Z(x)$, and substitute into the Bradley--Terry likelihood --- the intractable partition function cancels because it depends only on $x$ and appears in both terms of the difference. What is left is a purely supervised loss on preference pairs:
\[
\mathcal{L}_{\mathrm{DPO}} = -\log \sigma\!\left(\beta \log \frac{\pi_\theta(y_w \mid x)}{\pi_{\text{ref}}(y_w \mid x)} - \beta \log \frac{\pi_\theta(y_l \mid x)}{\pi_{\text{ref}}(y_l \mid x)}\right).
\]
The limits: it is offline, so the gradient never sees a sample the current policy would actually produce, and failure modes absent from the preference data go unfixed.

\textbf{Where it lives.} [NLP 8], ``Explain DPO. How does it relate to RLHF?'' and ``Derive the DPO loss from the RLHF objective.'' The RL mathematics underneath is [DL 17].

\textbf{The failure.} Presenting DPO as strictly better than PPO. The Staff answer is that DPO buys simplicity and loses on-policy exploration, which is why iterative DPO exists and why the frontier moved to online methods for anything with a verifiable reward.

\subsection*{21. GRPO and verifiable rewards}

\textbf{What you must be able to say.} GRPO is PPO with the critic deleted. For each prompt, sample $K$ responses, score them, and use the group's own statistics as the baseline: $\hat{A}_i = (r_i - \mathrm{mean}(r))/\mathrm{std}(r)$, assigned identically to every token of that response. That removes a full model copy with gradients and optimizer state (four resident models become three, or two with a rule-based verifier), and removes the weakest link --- a critic asked to predict per-token return from one sparse scalar at the end of a 10{,}000-token trace. It pairs with RLVR, where the reward is a program (unit tests, an answer checker) rather than a learned model, which shrinks the proxy gap and lets you apply far more optimization pressure. The known pathology: responses grow unboundedly long while benchmark scores stall --- a length-bias artifact of the normalization, not a capability gain.

\textbf{Where it lives.} [DL 17], ``Explain GRPO. Why did it replace PPO's critic for reasoning RL, and what breaks when you use it naively?'' and the runaway-length debugging question in the same chapter.

\textbf{The failure.} Believing verifiable rewards repeal Goodhart's law. The verifier is now the reward, so test-hacking is now the reward hacking --- hard-coded expected outputs, special-cased visible tests --- and the defense moves from KL budgets to test-suite quality.

\subsection*{22. In-context learning, chain of thought, and test-time compute}

\textbf{What you must be able to say.} In-context learning is not gradient learning: the examples condition the forward pass, and much of the gain comes from specifying the \emph{format} and the label space rather than the mapping --- which is why even partially random labels still help. Chain of thought works because it converts a problem the model must solve in one forward pass into a sequence of forward passes: it buys serial computation, not knowledge. It therefore helps on multi-step reasoning and does nothing (or hurts) on tasks solvable in one step. The Staff caveat: the stated chain is not a faithful trace of the computation, so it is not an explanation you can audit. Test-time compute is the same trade at the system level --- more samples or longer thinking per query instead of a larger model --- and it stops being cheaper when query volume is high, when the task lacks a verifier to select among samples, or when the budget exceeds the gap to the next model size.

\textbf{Where it lives.} [NLP 7], ``Why does chain-of-thought prompting improve reasoning? What are its limitations?'' and ``Explain in-context learning.'' The economics are [DL 17], ``When is test-time compute cheaper than a bigger model --- and when does it stop working?''

\textbf{The failure.} Treating the chain of thought as an explanation of the model's reasoning, and proposing ``add CoT'' as a general quality lever without noting that it multiplies latency and token cost on every request, including the ones that never needed it.

% ============================================================================
\section*{IV. Inference and Systems}
\addcontentsline{toc}{section}{IV. Inference and Systems}
% ============================================================================

\subsection*{23. The roofline: compute-bound versus bandwidth-bound}

\textbf{What you must be able to say.} Runtime is at least $\max(\mathrm{FLOPs}/P,\ \mathrm{bytes}/B)$, so every kernel is classified by its arithmetic intensity (FLOPs per byte moved) against the hardware ridge point $P/B$ --- about 295 FLOP/byte on an H100 (989 TFLOP/s BF16 over 3.35\,TB/s), about 155 on an A100. Large GEMMs sit far above it; elementwise ops, normalizations, and batch-1 decode sit far below at intensity near 1. Worked: decoding a 70B model in FP16 must read 140\,GB of weights per token, which at 2\,TB/s is a 70\,ms floor with the compute finishing in under a millisecond --- 99\% memory-stalled. This one lens explains quantization (fewer bytes), batching (amortize the weight read), GQA and paged KV (shrink KV traffic), and kernel fusion including FlashAttention, which reduces no FLOPs at all and only deletes HBM round trips by tiling the score matrix so the $n \times n$ array is never materialized. When a training job reports 20--25\% MFU against a healthy 40--50\%, look for launch gaps and small kernels, unoverlapped communication, dataloader stalls, pipeline bubbles, and checkpointing overhead --- and check tensor-core tile alignment, since a hidden size of 4100 instead of 4096 can cost 35\% of throughput.

\textbf{Where it lives.} [DL 26], ``Is it compute-bound or bandwidth-bound on an H100? Show your method,'' ``Your training job shows the GPU at 20\% utilization. Diagnose it,'' and the kernel-fusion question that uses FlashAttention as its example. The FlashAttention algorithm itself is [DL 5].

\textbf{The failure.} Reasoning about inference in FLOPs. Almost every serving optimization that matters attacks bytes, and a candidate who cannot classify an op in thirty seconds cannot justify any of them.

\subsection*{24. KV-cache arithmetic and the GQA/MLA ladder}

\textbf{What you must be able to say.} Under causal masking the keys and values of past tokens never change, so you cache them and each new token costs $O(n)$ attention instead of $O(n^2)$ recompute. The bytes are $2 \times n_{\text{layers}} \times n_{\text{kv heads}} \times d_{\text{head}} \times \text{seq} \times \text{bytes}$. Worked for Llama-70B with GQA (80 layers, 8 KV heads, $d_{\text{head}} = 128$, FP16): $2 \times 80 \times 8 \times 128 \times 2 \approx 0.33$\,MB per token, so 128K context is about 43\,GB for a single sequence. That number is why the ladder exists: MQA shares one KV head across all query heads, GQA groups them (Llama-2: 32 query heads, 8 KV heads, a $4\times$ cut), and MLA keeps all heads but caches a low-rank latent instead of per-head K and V. MLA needs a separate decoupled RoPE key because the position rotation sits between the two up-projections and does not commute with the low-rank reconstruction, so position cannot be applied to the compressed latent after the fact.

\textbf{Where it lives.} [DL 5], ``Calculate the KV cache memory for serving LLaMA 70B at 128K context length'' and ``Compare MHA, GQA, MQA, and MLA for KV-cache efficiency. Why does MLA need a decoupled RoPE key?'' The application-layer version is [NLP 12].

\textbf{The failure.} Forgetting the factor of 2 for K and V, or using query heads instead of KV heads and overstating the cache by $4$--$8\times$. Conceptually: not connecting cache size to the serving batch ceiling, which is the reason anyone asks.

\subsection*{25. Speculative decoding}

\textbf{What you must be able to say.} A small draft model proposes $k$ tokens; the target model verifies all of them in one forward pass; a modified rejection-sampling rule accepts a prefix and resamples at the first rejection, which makes the output distribution provably identical to sampling from the target. It wins because decode at low batch is bandwidth-bound, so verifying $k$ tokens costs barely more than generating one --- the extra FLOPs are free. Acceptance rates of 70--90\% give roughly $1.5$--$2\times$; DeepSeek-V3's multi-token-prediction head reports 85--90\% acceptance for about $1.8\times$. It stops helping when the draft is poorly aligned (low acceptance means wasted drafts), and when batch size is already large enough that decode has become compute-bound --- then you are spending real FLOPs for nothing.

\textbf{Where it lives.} [DL 19], ``Explain speculative decoding. When would you use it and when would you not?'' and ``Your speculative decoding setup is slower than standard decoding. What went wrong?'' The distributional-equivalence argument is also [NLP 12].

\textbf{The failure.} Saying it is an approximation that trades quality for speed. It is exact, and the acceptance rule is why --- a candidate who cannot say that has not understood the mechanism. The second failure is not knowing it degrades at high batch.

\subsection*{26. The quantization ladder}

\textbf{What you must be able to say.} FP16/BF16 is 2 bytes per parameter, INT8 one, INT4 half --- so a 70B model drops from 140\,GB to 35\,GB at INT4, which is what turns an eight-GPU deployment into a two-GPU one. The decision rule: FP8 where hardware supports it natively and you want throughput with minimal quality risk; weight-only INT8 as the safe default; INT4 (GPTQ, AWQ) when memory is the binding constraint and you have accepted a real quality cost. The mechanism behind the cost is activation outliers --- a few channels with much larger magnitudes dominate the quantization range, which is why AWQ protects salient channels and SmoothQuant migrates the difficulty into the weights. And the failure to anticipate: aggregate benchmark scores barely move while code, math, and long-tail multilingual degrade sharply, because those tasks depend on precise low-probability distinctions that rounding destroys.

\textbf{Where it lives.} [DL 14], ``INT8 vs INT4 vs FP8 quantization --- what is your decision framework?'' and ``Your quantized model performs well on benchmarks but fails on specific domains (code, math). Why?'' The memory-savings arithmetic is [DL 19]; the method comparison (GPTQ, AWQ, GGUF) is in both [DL 19] and [NLP 12].

\textbf{The failure.} Quoting an average benchmark delta as evidence quantization is safe. The Staff answer names the evaluation you would actually run --- domain-specific, long-context, and long-tail --- because that is where the loss hides.

\subsection*{27. Designing the LLM serving stack}

\textbf{What you must be able to say.} Size before architecting. A 70B model at INT4 is about 35\,GB of weights; on an 8$\times$H100 node with TP 8 that leaves roughly 600\,GB for KV cache, and at about 1.3\,GB per request at 4K context that is around 400 concurrent requests per node. With 200 output tokens at roughly 25\,ms per token, each request occupies about 5 seconds, so Little's law says 1000 QPS needs about 5000 concurrent requests --- about 13 replica nodes. Then the components: continuous batching so finished sequences leave the batch immediately, PagedAttention so KV memory is not over-reserved (2--4$\times$ more concurrency), chunked prefill so long prompts do not stall decode, prefix caching for the shared system prompt, TP inside the node on NVLink, and autoscaling on queue depth. Finally, be honest about the SLO: 200 tokens at 25\,ms is 5 seconds, which misses a 2-second p99 --- so you cap output length, add speculative decoding, or move to FP8.

\textbf{Where it lives.} [DL 19], ``Design an LLM serving stack that handles 1,000 queries per second with a p99 latency of 2 seconds for a 70B parameter model,'' plus the reasoning-model capacity replan in the same chapter.

\textbf{The failure.} Naming vLLM and stopping. The question tests whether you can size a system from memory and throughput arithmetic; a design with no numbers in it fails regardless of how many correct component names it contains.

\subsection*{28. The p99 latency budget}

\textbf{What you must be able to say.} A 250\,ms p99 SLO is decomposed across the funnel, each stage engineered against its slice and enforced with its own timeout --- which converts a latency problem into a measurable quality-degradation problem. Two facts justify obsessing over p99 rather than the median. First, sessions sample the distribution many times: 20 interactions means an 18\% chance of at least one p99-grade response. Second, fan-out amplifies: a scatter-gather over 100 shards completes when the slowest responds, so a per-shard p99 becomes a $1 - 0.99^{100} \approx 63\%$ event at request level --- to hold request p99 you need per-shard p99.99. The p99/p50 ratio is the queueing-health diagnostic: rising ratio at flat p50 means saturation or one misbehaving shard, not a workload change. When you are over budget, you cut candidate counts and rerank depth before you cut model quality.

\textbf{Where it lives.} [SR 8], ``You own a search endpoint with a 250\,ms server-side p99 SLO. Walk me through the latency budget.'' The LLM-serving analogue (p99 ten times p50) is [DL 19].

\textbf{The failure.} Quoting average latency, or budgeting to p50 and treating the tail as noise. At scale the tail is the typical experience of your heaviest users and of every fan-out consumer downstream.

\subsection*{29. Prompt injection and the trust boundary}

\textbf{What you must be able to say.} Prompt injection is not a jailbreak: a jailbreak is the user attacking the model's policy, while injection is a \emph{third party} attacking the application through content the model reads --- a retrieved document, a web page, a tool result. The root cause is architectural: instructions and data share one channel, and no amount of prompt engineering separates them reliably. The defense is therefore a system design, not a prompt: treat all retrieved and tool-returned content as untrusted, give the agent least privilege over tools (read-only by default, human confirmation for irreversible or outbound actions), constrain outputs by schema rather than by instruction, and enforce authorization at the tool boundary where the model cannot argue with it. Detection classifiers and instruction hierarchies raise the cost; they do not close the hole.

\textbf{Where it lives.} [NLP 13], ``What is prompt injection? Design a defense system for a production LLM application.'' The same question in the LLM-capability framing is [NLP 7]; agent-specific safety for browsing and code execution is [DL 20].

\textbf{The failure.} Proposing ``add a system prompt telling it to ignore instructions in documents.'' That answer says the candidate has not internalized that the model has no mechanism to distinguish the two channels, and it is the single most common wrong answer to this question.

% ============================================================================
\section*{V. Retrieval, RAG, and Recommendation}
\addcontentsline{toc}{section}{V. Retrieval, RAG, and Recommendation}
% ============================================================================

\subsection*{30. BM25 over an inverted index}

\textbf{What you must be able to say.} Mechanically: the analyzer maps the query to terms; each term has a postings list of document IDs, delta-encoded and compressed; the query walks the lists document-at-a-time with a top-$k$ heap, and WAND or block-max WAND skips any block whose stored maximum contribution cannot enter the heap --- which is safe, exact pruning, not approximation. That is how a billion documents return top-10 in tens of milliseconds. The formula multiplies IDF by a saturating term-frequency factor: $k_1$ controls saturation (the TF term asymptotes at $k_1 + 1$, so the fifth occurrence adds almost nothing --- the structural defense against keyword stuffing that raw TF--IDF lacks), and $b$ controls length normalization (at $b = 1$ a double-length document needs double the frequency for the same score). Defaults $k_1 \approx 1.2$, $b \approx 0.75$ were tuned on news prose; short-field corpora like product titles want lower $b$.

\textbf{Where it lives.} [SR 2], ``Walk me through, mechanically, how BM25 over an inverted index returns the top-10 results from a billion documents in under 50\,ms'' and ``What do BM25's $k_1$ and $b$ actually control?'' The probabilistic derivation is [NLP 1].

\textbf{The failure.} Reciting the formula with no index mechanics, or the reverse. Also: dismissing BM25 as legacy. It still carries the majority of production search traffic, it is unbeatable on exact identifiers and rare terms, and it is the half of a hybrid system that dense retrieval cannot replace.

\subsection*{31. The funnel: bi-encoder versus cross-encoder}

\textbf{What you must be able to say.} A bi-encoder embeds query and document \emph{independently}, so document vectors are precomputed offline and query time is a lookup plus cheap arithmetic --- sublinear in corpus size. A cross-encoder feeds query and document together through one model, so nothing precomputes and scoring is one model inference per candidate: at $10^8$ documents and tens of milliseconds each, that is years of GPU time per query. That asymmetry \emph{is} the funnel. Each stage spends more per candidate on fewer candidates: microseconds across millions at retrieval, fractions of a millisecond across thousands at L1, tens of milliseconds across hundreds at the cross-encoder. The quality gap is real (a cross-encoder models term interaction directly), which is why you buy it exactly where it can matter --- on candidates the cheap stages already vouched for.

\textbf{Where it lives.} [SR 4], ``Bi-encoder vs.\ cross-encoder: why not cross-encode everything, and what exactly does the cross-encoder buy?'' The funnel-as-Pareto-construction argument is [SR 1]; the architectural fundamentals are [DL 7]; the RAG framing is [NLP 9].

\textbf{The failure.} Explaining the quality difference without the decomposability argument. The reason you cannot cross-encode everything is not that it is expensive --- it is that nothing can be precomputed, which changes the complexity class, and a candidate who says only ``it's slower'' has missed the structural point.

\subsection*{32. Two-tower training: in-batch negatives and logQ}

\textbf{What you must be able to say.} Train with a softmax over in-batch negatives: every other document in the batch is a negative for this query, so a batch of $B$ gives $B - 1$ negatives at no extra encoding cost --- which is why production recipes use batches of 512--8192 with cross-device all-gather of embeddings (cheap: only the vectors move) and gradient caching when memory binds. But this is a \emph{sampled} softmax whose negatives are drawn from the training distribution, roughly proportional to popularity $Q(d)$, so popular items appear as negatives far too often and the model learns to suppress them. The logQ correction restores an unbiased estimate by subtracting the log sampling probability from each logit at training time, $s'(q,d) = s(q,d) - \log Q(d)$, with $Q$ estimated from stream frequencies; you serve with the raw score. Hard negatives sharpen the decision boundary but risk false negatives (unlabeled relevant documents), so they are mined and denoised, not just mined.

\textbf{Where it lives.} [SR 4], ``Why do bi-encoders train with such large batches? Explain in-batch negatives and the logQ correction.'' The recsys version is [SR 6]; the metric-learning objectives themselves are [DL 7] and [DL 3].

\textbf{The failure.} Confusing the logQ correction (a train-time logit adjustment that removes sampling bias) with serving-time popularity boosting. A close second: proposing hard negatives without mentioning the false-negative trap, which is how a hard-negative pipeline quietly teaches the model that correct answers are wrong.

\subsection*{33. HNSW and the ANN recall knobs}

\textbf{What you must be able to say.} HNSW is a navigable small-world proximity graph with a layer hierarchy: node levels are drawn geometrically so upper layers are sparse express networks and layer 0 holds everything. Search greedy-descends the upper layers to reach the right region in few hops, then runs a best-first beam search at layer 0 with a priority queue of size \texttt{efSearch}, stopping when the closest unexplored candidate is worse than the worst kept result. The beam is bounded backtracking, and \texttt{efSearch} is \emph{the} query-time recall knob --- latency grows near-linearly with it, recall with diminishing returns. \texttt{efConstruction} and $M$ set graph quality at build time; skimping there produces an index whose recall plateaus no matter how you tune search. Memory is vectors plus graph (about $2M \times 4$ bytes per node at layer 0, roughly 128\,B at $M = 16$), and it wants everything in RAM because traversal is random access.

\textbf{Where it lives.} [SR 3], ``How does HNSW work, and why is it fast?'' plus the recall-after-reindex and metadata-filter debugging questions in the same chapter. The one-page summary for a DL-side interview is [DL 7].

\textbf{The failure.} Calling it ``a graph index, it's approximate'' with no search procedure and no named knob. The follow-up is always about recall --- and the two things you must be able to say are that \texttt{efSearch} trades latency for recall at query time and that a selective metadata filter can collapse recall by disconnecting the graph.

\subsection*{34. Hybrid retrieval and RRF}

\textbf{What you must be able to say.} You cannot combine BM25 and cosine with a weighted sum, because the scores live on incompatible, \emph{query-dependent} scales: BM25 is unbounded and swells with query length and term rarity, cosine is bounded and concentrated in a narrow band, so any fixed $\alpha$ tuned on one query mix breaks on another. Per-query min-max or z-score normalization does not rescue it --- you would be normalizing by order statistics of about 100 candidates (noisy, outlier-dominated) and erasing absolute confidence, so a query where retrieval found nothing normalizes to look confident. The robust first ship is Reciprocal Rank Fusion, $\mathrm{RRF}(d) = \sum_i 1/(k + \mathrm{rank}_i(d))$ with $k \approx 60$: rank-based, hence invariant to any monotone rescaling, and $k$ sets how much cross-source agreement outweighs a single source's top hit. The end state is to stop fusing and feed both scores and ranks to the L1/L2 ranker as features, so the weights become learned and query-conditional.

\textbf{Where it lives.} [SR 4], ``You run BM25 and dense retrieval in parallel. How do you combine them --- and why does $0.5\cdot\mathrm{BM25} + 0.5\cdot\mathrm{cosine}$ fail?'' The lexical-versus-dense split argument is [SR 2]; the RAG framing is [NLP 9].

\textbf{The failure.} Knowing the RRF formula but not why raw fusion fails --- that scores as memorization. Proposing normalization as the fix is the trap's second layer, not the answer.

\subsection*{35. NDCG from first principles}

\textbf{What you must be able to say.} Build it in three moves, naming what each fixes. Cumulative gain sums the gains but ignores order. DCG divides the gain at rank $i$ by $\log_2(i + 1)$ to model decaying attention, but its scale depends on how many good documents the query has. NDCG divides by the ideal DCG (judged documents sorted by grade) so it lands in $[0,1]$ and averages across queries. Do the arithmetic without hesitation: grades $(3,0,2)$ under exponential gain $2^{\mathrm{rel}} - 1$ give $(7,0,3)$; discounts are $1$, $0.631$, $0.5$; so $\mathrm{DCG@}3 = 8.5$, and with ideal grades $(3,2,0)$, $\mathrm{IDCG@}3 = 7 + 3(0.631) = 8.893$, giving $\mathrm{NDCG@}3 \approx 0.956$. State your conventions, because both are choices: exponential gain concentrates learning on the top grades, and $\log_2(i+1)$ leaves rank 1 undiscounted.

\textbf{Where it lives.} [SR 7], ``Derive NDCG from first principles, then compute NDCG@3 for a ranking with grades $(3,0,2)$.'' The one-line catalog entries for NDCG, MRR, and MAP are [DL 23].

\textbf{The failure.} Writing the discount as $\log_2 i$ (which divides rank 1 by zero) --- an immediate sign the formula was memorized. Also: normalizing by the DCG of the shown list re-sorted rather than by the judged ideal, and treating unjudged documents as irrelevant without saying so.

\subsection*{36. Position bias and counterfactual LTR}

\textbf{What you must be able to say.} Clicks are missing-not-at-random: $P(\text{click} \mid d, p) = \theta_p \cdot P(\text{relevant} \mid d)$, where $\theta_p$ is the probability the user examined position $p$. Train on raw clicks and you learn to reproduce the ranker that produced the log, so incumbents at position 1 stay there and better new items never rise. The fix is inverse propensity scoring: weight each click by $1/\theta_p$ so the objective is unbiased with respect to examination. Propensities come from position randomization (a small swap-experiment budget), intervention harvesting from natural ranking variation, or EM over the logs. The cost is variance --- small propensities produce enormous weights, so you clip them. Raising position bias unprompted in any discussion of click-trained ranking is close to a pass/fail gate at Staff level.

\textbf{Where it lives.} [SR 5], ``Your click-trained ranker keeps favoring the items that have historically sat at position 1 --- better new items never rise'' and the propensity-estimation question in the same chapter. The counterfactual-estimator statistics are [SR 7].

\textbf{The failure.} Proposing to ``add more training data'' or ``retrain more often,'' both of which entrench the loop harder. The second failure is naming IPS without the variance consequence, which is where the follow-up always goes.

\subsection*{37. Cold start}

\textbf{What you must be able to say.} Split it: new item and new user are different problems. New items have no interaction history, so you rank them by content --- text, image, and category embeddings projected into the same space as the collaborative embeddings, so a brand-new item is scorable on day zero. New users get a popularity or segment prior, plus fast onboarding signals (first session behavior, context, device, locale), with the model updated within the session rather than at the next training run. The exploration piece is what separates a Staff answer: without deliberate exposure (an epsilon of traffic, Thompson sampling, or a boosted score for under-exposed items) a new item never accumulates the feedback that would let it rank, which is the same feedback loop as position bias. And say the metric consequence: cold-start quality does not show up in aggregate offline metrics, so it needs its own slice.

\textbf{Where it lives.} [SR 6], ``How do you handle the cold start problem for new users and new items?'' The bridge version in Volume I is [DL 12].

\textbf{The failure.} A list of techniques with no exploration and no measurement. The insight the question is probing is that cold start is a feedback-loop problem, not a feature-availability problem --- content features get the item scored, but only exposure gets it learned.

\subsection*{38. Calibration in ranking and ads}

\textbf{What you must be able to say.} Ranking metrics are invariant to any monotone transform of the scores, so a model can have AUC 0.95 and ECE 0.15 at the same time --- perfect ordering with badly wrong probabilities. That is fine if the score is only used to sort and fatal if it is used to \emph{decide}: ad rank is $\mathrm{pCTR} \times \mathrm{bid}$ and pricing follows from pCTR, budget pacing estimates spend from pCTR, and threshold rules (``suppress below 0.001'') shift when the scale shifts. Measure with expected calibration error, $\mathrm{ECE} = \sum_b \frac{|B_b|}{N}\left|\mathrm{acc}(B_b) - \mathrm{conf}(B_b)\right|$, plus a reliability diagram. Fix with temperature scaling, Platt scaling, or isotonic regression fit on held-out data --- none of which change the ranking. And know the systematic offender: negative downsampling shifts the base rate, so the predictions must be corrected back before they are used as probabilities.

\textbf{Where it lives.} [SR 6], ``Why do recommendation models need calibration? What happens if pCTR is systematically over-confident?'' The ECE definition, reliability diagrams, and the AUC-0.95/ECE-0.15 question are [DL 23].

\textbf{The failure.} Treating calibration as a nicety. The strong answer names a downstream consumer of the probability --- an auction, a pacing controller, an absolute threshold --- because that is what turns miscalibration from an aesthetic complaint into a revenue bug.

\subsection*{39. RAG: chunking, grounding, and evaluation}

\textbf{What you must be able to say.} The pipeline is ingest, chunk, embed, index, retrieve, rerank, assemble the prompt, generate with citations. Chunking is the decision people skip: 256--512 tokens with 10--20\% overlap is the working default, small chunks retrieve precisely but starve the generator of context, large chunks do the reverse, and the size should be tuned on a held-out set rather than inherited. Evaluate the two halves separately or you cannot debug: retrieval with recall@$k$ and NDCG against a judged set, generation with faithfulness (is every claim supported by the retrieved context?) and answer correctness. That decomposition is also the diagnosis when the system returns the right documents and still answers wrongly --- the failure is in the generator, the prompt assembly, or the ordering, not in retrieval, and the usual causes are context lost in the middle of a long prompt, conflicting passages, or a model that prefers its parametric memory to the provided context.

\textbf{Where it lives.} [NLP 9], ``Design a RAG system for a customer support chatbot,'' ``What chunking strategies exist for RAG? How do you choose chunk size?'', ``How would you evaluate a RAG system end-to-end?'', and ``Your RAG system is returning correct documents but the LLM's answers are still wrong.'' The retrieval stage in production is [SR 8]; the retrieval-versus-long-context decision is [DL 21].

\textbf{The failure.} Reporting one end-to-end quality number. Without separating retrieval from generation you cannot say which half regressed, and every real RAG debugging conversation starts exactly there.

% ============================================================================
\section*{VI. Classical ML, Evaluation, and Experimentation}
\addcontentsline{toc}{section}{VI. Classical ML, Evaluation, and Experimentation}
% ============================================================================

\subsection*{40. The XGBoost derivation}

\textbf{What you must be able to say.} Taylor-expand the round-$m$ loss to second order in the new tree's output, keeping $g_i$ and $h_i$ as the first and second derivatives at the current prediction, and add the regularizer $\gamma T + \frac{\lambda}{2}\sum_j w_j^2$. Grouping by leaf turns the objective into $T$ independent one-dimensional quadratics, $\sum_j [G_j w_j + \frac{1}{2}(H_j + \lambda)w_j^2] + \gamma T$, minimized at
\[
w_j^{*} = -\frac{G_j}{H_j + \lambda},
\qquad
\mathrm{Gain} = \frac{1}{2}\left[\frac{G_L^2}{H_L + \lambda} + \frac{G_R^2}{H_R + \lambda} - \frac{(G_L + G_R)^2}{H_L + H_R + \lambda}\right] - \gamma .
\]
Sanity-check it aloud: for squared loss ($g_i = \hat{y}_i - y_i$, $h_i = 1$) the leaf weight is the ridge-shrunk mean residual, recovering classic residual-fitting at $\lambda = 0$; for log loss it is a regularized Newton step in logit space. A split with negative gain is not worth its $\gamma$ --- cost-complexity pruning re-derived inside the boosting objective.

\textbf{Where it lives.} [CML 2], ``Derive the optimal leaf weight and the split-gain formula in XGBoost.'' The LambdaMART application of the same machinery is [SR 5].

\textbf{The failure.} Describing gradient boosting as ``fitting residuals.'' That is the squared-loss special case, not the definition --- the definition is functional gradient descent, and the interviewer asks this question precisely to find out which one you carry.

\subsection*{41. GBDT versus neural networks, and which GBDT}

\textbf{What you must be able to say.} On tabular data with heterogeneous, hand-engineered features, tuned GBDTs still match or beat MLPs, and the reasons are structural: axis-aligned splits are the right inductive bias for features with no metric relationship to each other, trees are invariant to monotone feature transforms and handle missing values natively, and they need far less tuning to get 95\% of the achievable quality. Neural networks win when the features become embeddings, sequences, high-cardinality IDs, or when multi-task and transfer matter. Between libraries: LightGBM's histogram plus leaf-wise growth is fastest on wide data; CatBoost's ordered target statistics handle very high-cardinality categoricals (a 40K-cardinality merchant ID) without the target leakage that naive target encoding introduces; XGBoost is the conservative default. And the standing caveat: a tree cannot extrapolate beyond the range of the training data, which bites on any feature with a trend.

\textbf{Where it lives.} [CML 2], ``Why does a gradient-boosted tree beat your MLP on this tabular dataset --- and when would it not?'' and ``XGBoost vs.\ LightGBM vs.\ CatBoost: 600K rows, 45 features of which 14 are categorical.'' The same decision in Volume I framing is [DL 24].

\textbf{The failure.} Answering with library preference rather than data properties. The question is always posed with a dataset description, and the strong answer maps each stated property (cardinality, row count, categorical share, latency budget) onto a specific mechanism.

\subsection*{42. Sizing an A/B test}

\textbf{What you must be able to say.} Convert to an absolute effect first: 1\% relative on a 2\% base is $\delta = 0.0002$. Then $n \approx 16\,p(1-p)/\delta^2$ per group at 80\% power and $\alpha = 0.05$ (the 16 is $2 \times 2.8^2$, from $1.96 + 0.84$). That gives $n \approx 16(0.02)(0.98)/4\times10^{-8} \approx 7.84$M per group, 15.7M total; at 2M eligible users per day, about eight days, rounded up to two full weeks so the test spans whole weekly cycles. Then the moves that separate levels: reframe as MDE rather than sample size (with the traffic we have, what is the smallest effect we can see? if that exceeds any plausible effect, do not run the test); note that only triggered users count, so a feature firing for 30\% of traffic triples the calendar time; and note that CTR is a ratio metric, so with user-level randomization the variance needs the delta method or a bootstrap over users, not over impressions.

\textbf{Where it lives.} [CML 6], ``We want to detect a 1\% relative lift on a 2\% click-through rate. How many users do we need, and how long should we run?'' The LLM-eval analogue is [DL 23].

\textbf{The failure.} Forgetting the relative-to-absolute conversion, or missing that $n$ scales with $\delta^{-2}$ --- halving the detectable effect costs four times the traffic. Both are where most candidates fail this question.

\subsection*{43. Sample ratio mismatch}

\textbf{What you must be able to say.} SRM means the observed allocation differs from the intended one --- 50/50 designed, 50.4/49.6 observed --- tested by chi-square against the intended ratio and flagged at $p < 0.001$. At experiment scale even a 0.2\% deviation is wildly improbable by chance, so an SRM means the assignment or logging pipeline is broken: asymmetric bot filtering, redirect latency dropping treatment users, a crash in one arm removing its worst-affected users, trigger logic evaluated differently across arms, or a dedup join on a field one arm populates differently. The consequence to state first and without hedging: an experiment with an SRM is \emph{not analyzable}. The arms are no longer comparable populations, and the users who dropped are usually the ones most affected. Debug it; never adjust for it.

\textbf{Where it lives.} [CML 6], ``Your treatment shows +2\% on the primary metric, but the SRM check fired. What now?''

\textbf{The failure.} Shipping anyway because the effect is large and the SRM is ``probably just logging,'' or proposing a statistical correction. This question exists to test exactly that instinct, and both answers fail it.

\subsection*{44. CUPED and variance reduction}

\textbf{What you must be able to say.} Sample size is linear in variance, so any variance you can explain \emph{without touching treatment} is traffic you do not have to buy. CUPED adjusts the metric as $Y' = Y - \theta(X - \bar{X})$ with $\theta^{*} = \mathrm{Cov}(Y,X)/\mathrm{Var}(X)$, where $X$ is a pre-experiment covariate --- almost always the same metric measured before the experiment. The expectation is unchanged and the variance is multiplied by $1 - \rho^2$: at $\rho = 0.7$ you need about 51\% of the users, so a two-week test becomes one week; at $\rho = 0.5$ the saving is 25\%; below $\rho \approx 0.3$ it is not worth the plumbing. The limits: it does nothing for new users (no pre-period), it is weak for rare conversions, and the covariate must be strictly pre-treatment --- a post-assignment covariate introduces real bias, which is the one way to get this genuinely wrong.

\textbf{Where it lives.} [CML 6], ``Explain CUPED to a PM, and quantify what it buys us.'' It also appears as the answer to the underpowered-test question in the same chapter.

\textbf{The failure.} Claiming CUPED increases power ``by reducing bias,'' or that it changes the estimand. It is ordinary regression adjustment chosen so the adjustment cannot itself be affected by treatment --- say that, and the quantification $1 - \rho^2$, and the question is over.

\subsection*{45. Interference and the marketplace test}

\textbf{What you must be able to say.} User-level randomization assumes one unit's outcome does not depend on another unit's assignment (SUTVA). On a two-sided marketplace that is false: treatment users who book more leave less supply for control users, so the measured gap overstates the true effect --- a pure transfer can look like growth. In social products the direction reverses: treatment users change their control-group friends' experience, diluting the effect toward zero. The designs that address it are cluster randomization (by region, community, or supply pool --- unbiased, but the effective sample size becomes the number of clusters, a brutal variance cost), switchback (region $\times$ time bucket with burn-in for carryover, well suited to pricing and dispatch), and budget-split designs in ads. A cheap diagnostic before any of that: vary the treatment allocation (10\% versus 50\%) and see whether the measured effect changes --- an allocation-dependent effect is direct evidence of interference.

\textbf{Where it lives.} [CML 6], ``Your marketplace A/B test shows +2\% GMV in treatment. Why might the true effect be smaller, or negative?''

\textbf{The failure.} Not spotting that a user-randomized marketplace test answers the wrong question at all. The interview signal is the recognition, not the vocabulary --- and the complete answer still runs the ordinary checks (SRM, novelty, whale concentration) because those explain many +2\% results too.

\subsection*{46. Offline metrics up, online flat or down}

\textbf{What you must be able to say.} Have five ordered causes ready. (1) The offline judgment set is biased by the logging policy --- it contains what the old ranker showed, so a model that surfaces genuinely new documents is scored against unjudged items counted as irrelevant. (2) Position and selection bias in click-derived labels, so offline gains partly measure agreement with the incumbent. (3) A funnel ceiling: a reranker cannot fix what retrieval never returned, so its offline gain on a fixed candidate set does not survive the real candidate distribution. (4) Metric mismatch --- NDCG on graded relevance is not CTR, and relevance is not revenue. (5) Systems effects: the new model added latency, and latency costs more engagement than the relevance gain returned. Then say what you would do: check the delta by query segment (head/torso/tail), the win/loss decomposition rather than the mean, and whether the online arm actually served what you evaluated.

\textbf{Where it lives.} [SR 7], ``Your reranker improved offline NDCG@10 by 3\% on the judgment set, but the A/B test shows flat CTR.'' The general-purpose version is [DL 23], ``Offline metrics went up but online metrics went down. Give 5 possible reasons''; the recsys version is [SR 6] and the query-understanding version is [SR 1].

\textbf{The failure.} Producing one cause and stopping, or jumping straight to ``the A/B test was underpowered.'' This is the most-asked debugging question in search and recommendation loops, and the expected shape is an ordered differential diagnosis with a cheap check for each branch.

% ============================================================================
\section*{VII. Craft}
\addcontentsline{toc}{section}{VII. Craft}
% ============================================================================

\subsection*{47. Implement causal multi-head attention}

\textbf{What you must be able to say.} You should be able to type this from memory in ten minutes: project to $Q$, $K$, $V$; reshape to (batch, heads, sequence, head\_dim) and transpose; scores $= QK^{\top}/\sqrt{d_{\text{head}}}$; add a causal mask of $-\infty$ above the diagonal \emph{before} the softmax; softmax over the last axis with the max subtracted for stability; multiply by $V$; transpose back, merge heads, apply the output projection. Say the shapes aloud as you go --- shape bugs are what actually fail this round. Know why the mask is additive and applied pre-softmax (a post-softmax mask leaves the normalizer wrong), and be ready to state the FLOP and memory cost of what you just wrote.

\textbf{Where it lives.} [DL 28], ``Implement multi-head self-attention with a causal mask in numpy.'' The runnable drills are in the program's \texttt{drills/} directory; the mechanism is [NLP 5] and the cost model is [DL 5].

\textbf{The failure.} Masking after the softmax, forgetting the $\sqrt{d_{\text{head}}}$ scale, or losing track of the head dimension in the reshape. Also: writing the code silently. In a coding round the narration is half the signal, and the candidate who says ``now I transpose so heads batch together'' reads very differently from one who does not.

\subsection*{48. Implement decoding: temperature, top-$k$, top-$p$, and a KV cache}

\textbf{What you must be able to say.} Sampling: divide logits by temperature ($T \to 0$ is greedy, $T > 1$ flattens); for top-$k$, keep the $k$ largest logits and set the rest to $-\infty$; for top-$p$, sort descending, take the cumulative softmax, and keep the smallest prefix whose mass exceeds $p$ --- always renormalizing after filtering, and always filtering on logits before the softmax rather than on probabilities after. The KV-cached greedy loop: run the full prompt once (prefill) and keep the per-layer K and V; then for each new token feed only the last token, append its K and V to the cache, and attend over the whole cache --- so the per-step cost is linear in context instead of quadratic. Handle the edge cases out loud: the first step has an empty cache, and the position index must come from the cache length, not from the input length.

\textbf{Where it lives.} [DL 28], ``Implement sampling from logits with temperature, top-$k$, and top-$p$'' and ``Implement greedy decoding with a KV cache.'' The decoding-strategy comparison is [NLP 10]; the serving consequences are [DL 19].

\textbf{The failure.} Applying top-$p$ to unsorted probabilities, forgetting to renormalize, or --- in the cache loop --- re-feeding the whole sequence each step, which silently makes the cache pointless. Interviewers check that last one specifically.

\subsection*{49. The impact-project deep dive}

\textbf{What you must be able to say.} Have one project rehearsed to survive forty minutes of drilling. The arc: the problem and why it mattered in business terms, the constraint that made it hard, two or three decisions with the alternative you rejected and \emph{why}, the result with a metric delta and a cost, and the reflection --- what you would do differently. Carry the numbers: scale, latency, cost, the size of the win, and the size of the error bar on the win. Expect the interviewer to pick one decision and go four levels down, so know the parts you did not build well enough to defend the interfaces, and be honest about what someone else did. The Staff signal is not that the project succeeded; it is that you can reconstruct the decision under questioning and show you were reasoning rather than following a recipe.

\textbf{Where it lives.} [DL 29], ``Walk me through your most impactful project. (Then 40 minutes of drilling.)'' The level calibration, the failure and conflict scenarios, and the research-taste questions are in the same chapter.

\textbf{The failure.} A narrative with no rejected alternatives and no numbers. ``We used a transformer and it worked'' answers nothing; the question is what else you considered, what you measured, and what it cost --- and the absence of a reflection reads as a candidate who has not learned from the work.

\subsection*{50. The deliberately vague design prompt}

\textbf{What you must be able to say.} Given a prompt like ``design memory for our assistant,'' the first ten minutes are not architecture. Ask what the product is and who the user is; establish scale (queries per day, sessions per user, retention horizon); establish the constraints that actually bind (latency, cost per query, privacy and deletion, multi-tenancy); and propose the success metric before proposing a system, because the metric determines the design. Then state your interpretation back --- ``so we are building X for Y, optimizing Z, under constraint W'' --- and get agreement before drawing anything. Sketch the simplest system that could work, name where it breaks first, and only then add machinery. Timebox: scope and metric by minute ten, a first architecture by minute twenty, then depth wherever the interviewer pushes.

\textbf{Where it lives.} [DL 29], ``I'm going to give you a deliberately vague prompt --- `design memory for our assistant.' Run the first ten minutes.'' The domain-specific design questions are [DL 19], [DL 22], [DL 24], [DL 27], [SR 4], [SR 6], and [NLP 9].

\textbf{The failure.} Drawing boxes in minute two. The prompt is vague on purpose: it tests whether you convert ambiguity into a scoped problem with a stated metric, which is the single clearest Staff-versus-Senior discriminator in the design round. The mirror-image failure is spending all thirty minutes gathering requirements and never committing to a design.

% ============================================================================
\section*{The Checklist}
\addcontentsline{toc}{section}{The Checklist}
% ============================================================================

Cover the entries above and work down this list. For each name, speak the substance out loud for ninety seconds without notes. Mark it green only if the number or the formula came out with it.

\begin{center}
\small
\begin{tabularx}{\textwidth}{@{}r X r X@{}}
\toprule
\multicolumn{2}{@{}l}{\textbf{I. Foundations and architectures}} & \multicolumn{2}{l@{}}{\textbf{III. Training and post-training}} \\
\midrule
1 & Softmax cross-entropy gradient & 13 & The memory ledger \\
2 & Perplexity, cross-entropy, bits & 14 & Parallelism: DP, TP, PP, ZeRO, EP \\
3 & Double descent & 15 & Loss spikes and NaNs \\
4 & 6ND and Chinchilla & 16 & The pretraining data pipeline \\
5 & Self-attention and $\sqrt{d_k}$ & 17 & LoRA: count, ratio, reason \\
6 & Attention versus FFN cost & 18 & Catastrophic forgetting \\
7 & Encoder, decoder, encoder--decoder & 19 & RLHF, the KL anchor, reward hacking \\
8 & Vanishing gradients and LSTMs & 20 & DPO \\
9 & Pre-norm, RMSNorm, SwiGLU & 21 & GRPO and verifiable rewards \\
10 & RoPE and context extension & 22 & ICL, chain of thought, test-time compute \\
11 & Mixture of experts & & \\
12 & Tokenization and the tokenizer tax & & \\
\bottomrule
\end{tabularx}

\medskip

\begin{tabularx}{\textwidth}{@{}r X r X@{}}
\toprule
\multicolumn{2}{@{}l}{\textbf{IV. Inference and systems}} & \multicolumn{2}{l@{}}{\textbf{V. Retrieval, RAG, recommendation}} \\
\midrule
23 & The roofline and MFU & 30 & BM25 over an inverted index \\
24 & KV-cache arithmetic, GQA, MLA & 31 & The funnel: bi- versus cross-encoder \\
25 & Speculative decoding & 32 & In-batch negatives and logQ \\
26 & The quantization ladder & 33 & HNSW and the recall knobs \\
27 & The LLM serving stack & 34 & Hybrid retrieval and RRF \\
28 & The p99 latency budget & 35 & NDCG from first principles \\
29 & Prompt injection & 36 & Position bias and counterfactual LTR \\
 & & 37 & Cold start \\
 & & 38 & Calibration in ranking and ads \\
 & & 39 & RAG: chunking, grounding, evaluation \\
\bottomrule
\end{tabularx}

\medskip

\begin{tabularx}{\textwidth}{@{}r X r X@{}}
\toprule
\multicolumn{2}{@{}l}{\textbf{VI. Classical ML and experimentation}} & \multicolumn{2}{l@{}}{\textbf{VII. Craft}} \\
\midrule
40 & The XGBoost derivation & 47 & Implement causal multi-head attention \\
41 & GBDT versus neural nets & 48 & Implement decoding and a KV cache \\
42 & Sizing an A/B test & 49 & The impact-project deep dive \\
43 & Sample ratio mismatch & 50 & The vague design prompt \\
44 & CUPED & & \\
45 & Interference in marketplaces & & \\
46 & Offline up, online flat & & \\
\bottomrule
\end{tabularx}
\end{center}

\begin{keyinsight}[The Only Scoring Rule That Matters]
Green is not ``I know this.'' Green is ``I said it out loud, in under ninety seconds, and the number was in it.'' Everything else is amber. On the last day before a loop, read nothing --- speak the ambers.
\end{keyinsight}
